% ID: SEL-D0046-OA-E03
% ESTADO: validada. S047.
Denotamos por $P$ el suceso «sale un número par». Hay tres pares en A y dos en B. La moneda asigna probabilidad $1/2$ a cada urna.

\textbf{a)} Por la regla del producto,
$P(A\cap P)=P(A)P(P\mid A)=\frac12\cdot\frac37=\boxed{\frac{3}{14}}$.

\textbf{b)} Sumamos los dos caminos disjuntos que llevan a un número par:
\[
P(P)=\tfrac12\cdot\tfrac37+\tfrac12\cdot\tfrac25=\boxed{\tfrac{29}{70}}.
\]
\textbf{c)} Al saber que el número es par, restringimos el espacio muestral a ese suceso. Por Bayes,
$P(A\mid P)=\frac{3/14}{29/70}=\boxed{\frac{15}{29}}$.
